\chapter*{Содержание}
\pdfbookmark[0]{Содержание}{contents}
\begingroup
\makeatletter\@starttoc{toc}\makeatother
\endgroup

\vspace{1.1em}
\section*{С чего начать}
\textbf{Для первого чтения} достаточно знать, как умножить матрицу на вектор,
что такое производная и как устроено комплексное число. Найдите на рисунке в главе~\ref{primer:chapter} треугольник
эталонной сетки GMsFEM и восстановите поле в его вершине.
В главе~\ref{walk:chapter} можно повторить этот прогноз в Python.
В разделе~\ref{code:operator} (с.\,\pageref{code:operator}) разберите, как
из матриц слоя получается оператор над коэффициентами поля.

\textbf{Для практикума} пройдите путь от расчётного поля до прогноза
в главах~\ref{math:fem}--\ref{walk:chapter}. Затем попробуйте объяснить
по таблице в главе~\ref{exp:chapter}, сколько точности теряется при сжатии
и сколько добавляет сеть.

\textbf{Для углублённого чтения} обратитесь к главе~\ref{op:chapter}
за условиями устойчивости и оценками ошибки. В главе~\ref{adv:chapter}
показано, что меняется при переходе к комплексному базису и нескольким ветвям сети.
